\documentclass[12pt, twoside]{article}
\input{../preamble}

\fancyhead[L]{La Scuola d'Italia / Dr.\ Huson / IB Math (b) \\* 14 March 2026}
\fancyhead[R]{\fbox{\begin{minipage}{6cm}\footnotesize Student ID:\\[0.5cm]\end{minipage}}}
\fancyfoot[C]{\thepage}

\begin{document}
\subsubsection*{7.11 PreQuiz: Logarithmic and Exponential Functions
(calculator)}

\begin{enumerate}[itemsep=0.15cm]

%--- Problem 1: Quarterly compound interest [6 marks] ---
%--- Source: Original, modelled on 7-10PreQuiz_Log-functions problem 4 ---
\item[1.] [Maximum mark: 6]

Lina invests \$2400 in an account that pays interest of $6\,\%$ per
annum, compounded quarterly. No further deposits or withdrawals are made.

\begin{enumerate}[itemsep=0.1cm]
  \item Write down the quarterly growth factor for the value of the
  account. \hfill [1]
  \item Find the value of the account at the end of $5$ years, correct
  to the nearest cent. \hfill [2]
  \item Find the smallest number of full years required for the value of
  the account to first exceed \$3400. \hfill [3]
\end{enumerate}

\begin{tikzpicture}
  \draw (-0.6,0) rectangle (11,11);
  \foreach \y in {1,2,...,10} \draw[dotted] (0,\y)--(10.5,\y);
\end{tikzpicture}

\newpage

%--- Problem 2: Continuous exponential growth [6 marks] ---
%--- Source: Original, modelled on 7-10PreQuiz_Log-functions problem 3 ---
\item[2.] [Maximum mark: 6]

The number of members of an online study group is modelled by
$$M(t) = 5000 e^{kt}, \qquad k > 0$$
where $t$ is the number of years after the group is created. After
$2$ years there are $7200$ members.

\begin{enumerate}[itemsep=0.1cm]
  \item Find the value of $k$. \hfill [2]
  \item Find the number of members after $6$ years, correct to the
  nearest whole number. \hfill [2]
  \item Find the time, in years, when the number of members first
  reaches $10\,000$. Give your answer correct to two decimal places.
  \hfill [2]
\end{enumerate}

\begin{tikzpicture}
  \draw (-0.6,0) rectangle (11,12);
  \foreach \y in {1,2,...,11} \draw[dotted] (0,\y)--(10.5,\y);
\end{tikzpicture}

\end{enumerate}
\end{document}
